\subsection{Distance in the Plane} 
The \emph{distance} between $A$ and $B$, written $d(A,B)$, is the length of $\overline{AB}$, written $AB$. It is easy to find for \makeblank[1.5in] and \makeblank[1.5in] segments.
\begin{Example}
Plot the points $A=(-3,2)$, $B=(1,2)$, $C=(2,2)$, $D=(2,-1)$. Then find $AB$ and $CD$.
\begin{lpic}{}{5}
\end{lpic}
\end{Example}
In general:
\begin{itemize}
\item If $A=(x,y_A)$ and $B=(x,y_B)$ (same \makeblanksmall coordinate), $AB=\makeblank[1in]$
\item If $A=(x_A,y)$ and $B=(x_B,y)$ (same \makeblanksmall coordinate), $AB=\makeblank[1in]$
\end{itemize}
For other segments, we use the \makeblank[1in] theorem.
\begin{theorem}[Pythagorean Theorem]
If the legs of a right triangle have lengths $a$ and $b$, and the hypotenuse has length $c$, then
\[
a^2+b^2=c^2
\]
\end{theorem}
\begin{Example}
Plot the points $A=(-3,-2)$ and $B=(1,1)$. Then find $AB$.
\begin{lpic}{}{5}
\end{lpic}
\end{Example}
\begin{displaybox}[Distance Formula]
If $A=(x_A,y_A)$ and $B=(x_B,y_B)$, then
\[
d(A,B)=\sqrt{\left(x_A-x_B\right)^2+\left(y_A-y_B\right)^2}
\]
\end{displaybox}